% !TEX root = ../main.tex
%% Corrigés — Chapitre 10 — Polynômes et fractions rationnelles
%% Résultats vérifiés par verifications/chapitre_10.py.

\corrigesChapitre{10}{Polynômes et fractions rationnelles}
\corrigesSouples

\begin{blocLivre}{cor-10-1}
\begin{corrige}{1}{Degrés}
\demarche{Le degré d'un produit s'ajoute ; celui d'une somme peut chuter.}
\begin{questions}
  \item $\deg=5$, coefficient dominant $2$ ;\quad $\deg=6$, dominant $1$ ;\quad
    $(X^{3}+1)-(X^{3}-X)=X+1$ : degré $1$, dominant $1$.
  \item $P=X^{2}+X$ et $Q=-X^{2}$ : $P+Q=X$, de degré $1<2$.
  \item $\deg(P\circ Q)=\deg P\times\deg Q=\res{6}$.
\end{questions}
\end{corrige}
\end{blocLivre}

\begin{blocLivre}{cor-10-2}
\begin{corrige}{2}{Division euclidienne}
\demarche{Éliminer le terme dominant à chaque étape — comme pour les entiers.}
\begin{questions}
  \item $X^{3}+2X^{2}-X+1=(X^{2}+1)(X+2)+(-2X-1)$ :
    quotient $\res{X+2}$, reste $\res{-2X-1}$ (de degré $<2$).
  \item $X^{4}-1=(X-2)(X^{3}+2X^{2}+4X+8)+15$ :
    quotient $\res{X^{3}+2X^{2}+4X+8}$, reste $\res{15}$.
  \item Le reste de la division par $X-a$ est la constante $P(a)$
    (écrire $P=(X-a)Q+R$ et évaluer en $a$). Ici
    $P(2)=2^{4}-1=15$ : conforme.
\end{questions}
\end{corrige}
\end{blocLivre}

\begin{blocLivre}{cor-10-3}
\begin{corrige}{3}{Racines évidentes et factorisation}
\demarche{Une racine, une division — puis on recommence.}
\begin{questions}
  \item $P(1)=1-6+11-6=0$, et la division donne
    $P=(X-1)(X^{2}-5X+6)$.
  \item $X^{2}-5X+6=(X-2)(X-3)$, d'où
    $\res{P=(X-1)(X-2)(X-3)}$.
  \item Somme des racines $1+2+3=6=-\dfrac{-6}{1}$ ; produit
    $6=(-1)^{3}\dfrac{-6}{1}$ : Viète est vérifié.
\end{questions}
\end{corrige}
\end{blocLivre}

\begin{blocLivre}{cor-10-4}
\begin{corrige}{4}{Multiplicité d'une racine}
\demarche{Les dérivées successives comptent la multiplicité.}
\begin{questions}
  \item $P(1)=0$ ; $P'=4X^{3}-15X^{2}+18X-7$, $P'(1)=0$ ;
    $P''=12X^{2}-30X+18$, $P''(1)=0$ ; $P'''=24X-30$,
    $P'''(1)=-6\neq0$.
  \item Trois dérivées nulles puis une non nulle : $1$ est racine de
    multiplicité $\res{3}$, donc $(X-1)^{3}$ divise $P$ ; le quotient, de
    degré $1$ et de produit des racines donné par $P(0)=2$, est $X-2$ :
    \[
      \res{P=(X-1)^{3}(X-2)}.
    \]
\end{questions}
\end{corrige}
\end{blocLivre}

\begin{blocLivre}{cor-10-5}
\begin{corrige}{5}{Relations de Viète}
\demarche{Somme et produit se lisent sur les coefficients extrêmes.}
\begin{questions}
  \item Racines de $X^{2}-7X+10$ : $\res{2\text{ et }5}$.
  \item $x_1+x_2+x_3=-\dfrac{-2}{1}=\res{2}$ et
    $x_1x_2x_3=(-1)^{3}\dfrac{-4}{1}=\res{4}$.
  \item $(x_1+x_2+x_3)^{2}=s_2+2(x_1x_2+x_1x_3+x_2x_3)$, et le
    coefficient de $X$ donne $x_1x_2+x_1x_3+x_2x_3=3$ : d'où
    $s_2=2^{2}-2\times3=\res{-2}$ — négatif, preuve que les racines ne
    sont pas toutes réelles !
\end{questions}
\end{corrige}
\end{blocLivre}

\begin{blocLivre}{cor-10-6}
\begin{corrige}{6}{L'algorithme de Horner}
\demarche{Une multiplication et une addition par coefficient, et la division par $X-c$ offerte.}
\begin{questions}
  \item Coefficients $2,-3,0,1,-5$ ; avec $c=2$ :
    $2\to2\times2-3=1\to1\times2+0=2\to2\times2+1=5\to5\times2-5=5$ :
    $\res{P(2)=5}$.
  \item Les valeurs intermédiaires sont les coefficients du quotient :
    $\res{P=(X-2)(2X^{3}+X^{2}+2X+5)+5}$.
  \item Quatre multiplications, contre $4+3+2+1=10$ pour l'évaluation
    naïve des puissances : Horner est linéaire en le degré.
\end{questions}
\end{corrige}
\end{blocLivre}

\begin{blocLivre}{cor-10-7}
\begin{corrige}{7}{Factoriser sur $\Cx$ puis sur $\R$}
\demarche{Racines complexes d'abord, regroupement des conjuguées ensuite.}
\begin{questions}
  \item $X^{4}-1=(X-1)(X+1)(X-\mathrm{i})(X+\mathrm{i})$ sur $\Cx$, et
    $\res{(X-1)(X+1)(X^{2}+1)}$ sur $\R$.
  \item $X^{3}-1=(X-1)(X-j)(X-j^{2})$ sur $\Cx$, et
    $\res{(X-1)(X^{2}+X+1)}$ sur $\R$.
  \item $X^{4}+X^{2}+1=(X^{2}+1)^{2}-X^{2}
    =\res{(X^{2}+X+1)(X^{2}-X+1)}$ sur $\R$ (les deux facteurs ont
    $\Delta=-3<0$) ; sur $\Cx$, chacun se scinde avec les racines
    $-j,-j^{2}$ et $j? $ — précisément $\pm j,\pm j^{2}$ selon le facteur.
\end{questions}
\end{corrige}
\end{blocLivre}

\begin{blocLivre}{cor-10-8}
\begin{corrige}{8}{PGCD de deux polynômes}
\demarche{Euclide, puis normaliser ; ou comparer les factorisations.}
\begin{questions}
  \item $X^{3}-3X+2=(X^{2}-1)\,X+(-2X+2)$, puis
    $X^{2}-1=(-2X+2)\bigl(-\tfrac12X-\tfrac12\bigr)+0$ : le dernier reste
    non nul est $-2X+2$, que l'on normalise :
    $\res{A\wedge B=X-1}$.
  \item $A=(X-1)^{2}(X+2)$ et $B=(X-1)(X+1)$ : le facteur commun est
    bien $X-1$.
\end{questions}
\end{corrige}
\end{blocLivre}

\begin{blocLivre}{cor-10-9}
\begin{corrige}{9}{La formule de Taylor polynomiale}
\demarche{Les dérivées en $a$ deviennent les coefficients en puissances de $X-a$.}
\begin{questions}
  \item Pour $P=X^{3}$ : $P(1)=1$, $P'(1)=3$, $P''(1)=6$, $P'''(1)=6$,
    d'où
    \[
      \res{X^{3}=1+3(X-1)+3(X-1)^{2}+(X-1)^{3}}.
    \]
  \item Le développement du membre de droite redonne $X^{3}$
    (c'est $(1+(X-1))^{3}$, par le binôme !).
  \item $1+3h+3h^{2}+h^{3}=(1+h)^{3}$ : la formule de Taylor en $1$
    n'est autre que le changement de variable $X=1+h$.
\end{questions}
\end{corrige}
\end{blocLivre}

\begin{blocLivre}{cor-10-10}
\begin{corrige}{10}{Racines doubles et paramètre}
\demarche{Racine double = racine commune de $P$ et $P'$ : le paramètre se lit sur $P'$.}
\begin{questions}
  \item C'est la caractérisation du cours ($m\geq2$ équivaut à
    $P(a)=P'(a)=0$).
  \item $P_\lambda'=3X^{2}-3$ s'annule en $\pm1$. Racine double en $1$ :
    $P_\lambda(1)=\lambda-2=0$, soit $\lambda=2$, et alors
    $\res{P_2=(X-1)^{2}(X+2)}$ ; racine double en $-1$ :
    $\lambda=-2$ et $\res{P_{-2}=(X+1)^{2}(X-2)}$.
  \item La fonction $x\mapsto x^{3}-3x$ a pour maximum local $2$ (en
    $-1$) et minimum local $-2$ (en $1$) : l'équation $x^{3}-3x=-\lambda$
    a trois solutions distinctes si et seulement si
    $\res{-2<\lambda<2}$.
\end{questions}
\end{corrige}
\end{blocLivre}

\begin{blocLivre}{cor-10-11}
\begin{corrige}{11}{La rigidité des polynômes}
\demarche{Trop de racines tue le polynôme.}
\begin{questions}
  \item Un polynôme non nul de degré $\leq n$ a au plus $n$ racines :
    s'il s'annule en $n+1$ points, il est nul.
  \item $P-X$ est de degré $\leq n$ et s'annule en les $n+1$ points
    $0,1,\dots,n$ : $P-X=0$, donc $\res{P=X}$.
  \item $Q=P(X+1)-P(X)$ s'annule en tout réel : $Q=0$, donc $P$ prend la
    même valeur sur tous les entiers... plus précisément : $P-P(0)$
    s'annule en $0,1,2,\dots$, une infinité de points : $P=P(0)$,
    \res{\text{constant}}.
\end{questions}
\end{corrige}
\end{blocLivre}

\begin{blocLivre}{cor-10-12}
\begin{corrige}{12}{Polynômes réels et racines conjuguées}
\demarche{Conjuguer une relation à coefficients réels : la conjugaison ne concerne que les racines non réelles.}
\begin{questions}
  \item En conjuguant $P(z)=0$ (les coefficients sont réels) :
    $P(\overline{z})=\overline{P(z)}=0$. L'argument s'applique aussi aux
    dérivées : $z$ et $\overline{z}$ annulent les mêmes
    $P^{(k)}$, donc ont la même multiplicité.
  \item $1-\mathrm{i}$ est racine car c'est la conjuguée de la racine
    \emph{non réelle} $1+\mathrm{i}$ ; mais $\sqrt2$ est réelle : sa
    « conjuguée » est elle-même, et rien n'impose $-\sqrt2$. En la
    choisissant :
    \[
      \res{(X^{2}-2X+2)(X^{2}-2)=X^{4}-2X^{3}+4X-4}.
    \]
  \item Chapitre 6 : TVI et limites infinies de signes opposés.
    Chapitre 10 : les racines non réelles vont par paires conjuguées — un
    degré impair ne peut s'épuiser en paires.
\end{questions}
\end{corrige}
\end{blocLivre}

\begin{blocLivre}{cor-10-13}
\begin{corrige}{13}{$X^{n}-1$ et les racines de l'unité}
\demarche{Apparier $\omega$ et $\overline{\omega}$ : chaque paire donne un trinôme réel.}
\begin{questions}
  \item $X^{n}-1=\displaystyle\prod_{k=0}^{n-1}\bigl(X-\mathrm{e}^{2\mathrm{i}k\pi/n}\bigr)$.
  \item $(X-\omega)(X-\overline{\omega})=X^{2}-2\cos\theta\,X+1$ pour
    $\omega=\mathrm{e}^{\mathrm{i}\theta}$, d'où
    \[
      X^{5}-1=(X-1)\Bigl(X^{2}-2\cos\tfrac{2\pi}{5}X+1\Bigr)
      \Bigl(X^{2}-2\cos\tfrac{4\pi}{5}X+1\Bigr).
    \]
  \item Avec $\cos\dfrac{2\pi}{5}=\dfrac{\sqrt5-1}{4}$ et
    $\cos\dfrac{4\pi}{5}=-\dfrac{1+\sqrt5}{4}$ (chapitre 3) :
    \[
      \res{X^{5}-1=(X-1)
      \Bigl(X^{2}-\tfrac{\sqrt5-1}{2}X+1\Bigr)
      \Bigl(X^{2}+\tfrac{\sqrt5+1}{2}X+1\Bigr)}.
    \]
\end{questions}
\end{corrige}
\end{blocLivre}

\begin{blocLivre}{cor-10-14}
\begin{corrige}{14}{Bézout chez les polynômes}
\demarche{Pas de racine commune dans $\Cx$, ou un reste constant : deux preuves de la primalité relative.}
\begin{questions}
  \item Les racines de $A$ sont $\pm\mathrm{i}$, celle de $B$ est $1$ :
    aucune racine commune, donc $A\wedge B=1$. Par Euclide :
    $X^{2}+1=(X-1)(X+1)+2$, reste constant non nul — même conclusion.
  \item De la division : $2=(X^{2}+1)-(X+1)(X-1)$, d'où
    \[
      \res{\tfrac12\,(X^{2}+1)-\tfrac12\,(X+1)\,(X-1)=1} :
    \]
    $U=\tfrac12$, $V=-\tfrac12(X+1)$.
  \item En multipliant par $C$ : $(UC)A+(VC)B=C$ pour tout $C$.
\end{questions}
\end{corrige}
\end{blocLivre}

\begin{blocLivre}{cor-10-15}
\begin{corrige}{15}{Interpolation de Lagrange}
\demarche{Les $L_i$ valent 1 en leur nœud, 0 ailleurs : la combinaison est immédiate.}
\begin{questions}
  \item $L_0=\dfrac{(X-1)(X-2)}{2}$, $L_1=-X(X-2)$,
    $L_2=\dfrac{X(X-1)}{2}$, et
    \[
      P=1\cdot L_0+3\cdot L_1+2\cdot L_2
      =\res{-\frac32X^{2}+\frac72X+1}.
    \]
    (Contrôle : $P(0)=1$, $P(1)=3$, $P(2)=2$.)
  \item $Q$ convient si et seulement si $Q-P$ s'annule en $0,1,2$ :
    $\res{Q=P+X(X-1)(X-2)\,R},\ R\in\R[X]$.
  \item L'unique interpolant de degré $\leq3$ des points $(k,k)$ est
    $P=X$ (exercice 11) : l'interpolation \emph{reconnaît} les polynômes
    de petit degré.
\end{questions}
\end{corrige}
\end{blocLivre}

\begin{blocLivre}{cor-10-16}
\begin{corrige}{16}{Fractions : forme irréductible et partie entière}
\demarche{Simplifier d'abord, diviser ensuite.}
\begin{questions}
  \item $X^{3}+1=(X+1)(X^{2}-X+1)$ et $X^{2}-1=(X+1)(X-1)$ :
    \[
      \res{F=\frac{X^{2}-X+1}{X-1}}.
    \]
    Zéros : les racines de $X^{2}-X+1$ (complexes, $-j$ et $-j^{2}$,
    simples) ; pôle : $1$, simple. ($-1$ n'est \emph{pas} un pôle : il a
    disparu dans la simplification.)
  \item $X^{2}-X+1=X(X-1)+1$ :
    $\res{F=X+\dfrac{1}{X-1}}$.
  \item Quand $x\to\pm\infty$, $F(x)-x=\dfrac{1}{x-1}\to0$ : la droite
    $\res{y=x}$ est asymptote oblique.
\end{questions}
\end{corrige}
\end{blocLivre}

\begin{blocLivre}{cor-10-17}
\begin{corrige}{17}{Éléments simples : pôles simples}
\demarche{Multiplier par $(X-\lambda)$ et évaluer en $\lambda$.}
\begin{questions}
  \item $\dfrac{1}{X^{2}-1}
    =\res{\dfrac12\cdot\dfrac{1}{X-1}-\dfrac12\cdot\dfrac{1}{X+1}}$.
  \item En $0$ : $\dfrac{0+2}{0+1}=2$ ; en $-1$ : $\dfrac{-1+2}{-1}=-1$ :
    $\dfrac{X+2}{X(X+1)}=\res{\dfrac{2}{X}-\dfrac{1}{X+1}}$.
  \item $\dfrac{1}{X(X^{2}+1)}=\res{\dfrac1X-\dfrac{X}{X^{2}+1}}$
    (le coefficient en $0$ vaut $1$, et l'on identifie le reste).
  \item Sur $\intOO{1}{+\infty}$ :
    $\displaystyle\int\frac{\mathrm{d}x}{x^{2}-1}
    =\res{\frac12\ln\frac{x-1}{x+1}}+C$.
\end{questions}
\end{corrige}
\end{blocLivre}

\begin{blocLivre}{cor-10-18}
\begin{corrige}{18}{La dérivée logarithmique}
\demarche{Dériver un produit, diviser par le produit : chaque facteur livre son terme.}
\begin{questions}
  \item Par récurrence sur le nombre de facteurs, ou directement : la
    dérivée de $P=c\prod(X-\lambda_i)^{m_i}$ est une somme où chaque terme
    remplace un facteur $(X-\lambda_i)^{m_i}$ par sa dérivée
    $m_i(X-\lambda_i)^{m_i-1}$ ; en divisant par $P$ :
    $\dfrac{P'}{P}=\sum\dfrac{m_i}{X-\lambda_i}$.
  \item Pour $P=(X-1)^{2}(X+2)$ :
    $\dfrac{P'}{P}=\res{\dfrac{2}{X-1}+\dfrac{1}{X+2}}$ — vérification
    directe en dérivant.
  \item Les multiplicités sont les coefficients de la décomposition de
    $P'/P$ : une racine est simple si et seulement si son coefficient
    vaut $1$.
\end{questions}
\end{corrige}
\end{blocLivre}

\begin{blocLivre}{cor-10-19}
\begin{corrige}{19}{Tchebychev : racines et factorisation}
\demarche{L'identité $T_n(\cos t)=\cos(nt)$ fournit toutes les racines d'un coup.}
\begin{questions}
  \item Par récurrence double : si $\deg T_{n+1}=n+1$ avec dominant
    $2^{n}$ et $\deg T_n=n$, alors $2X\,T_{n+1}$ domine, de degré $n+2$
    et de coefficient $2^{n+1}$.
  \item $T_n(x_k)=\cos\Bigl(n\cdot\dfrac{(2k+1)\pi}{2n}\Bigr)
    =\cos\Bigl(\dfrac{(2k+1)\pi}{2}\Bigr)=0$. Les $x_k$ sont deux à deux
    distincts : les $\dfrac{(2k+1)\pi}{2n}$ sont dans $\intOO{0}{\pi}$, où
    $\cos$ est injective.
  \item Voilà $n$ racines distinctes pour un polynôme de degré $n$ :
    toutes simples, toutes dans $\intOO{-1}{1}$, et
    \[
      \res{T_n=2^{n-1}\prod_{k=0}^{n-1}
      \Bigl(X-\cos\tfrac{(2k+1)\pi}{2n}\Bigr)}.
    \]
\end{questions}
\end{corrige}
\end{blocLivre}

\begin{blocLivre}{cor-10-20}
\begin{corrige}{20}{Sommes de Newton}
\demarche{Chaque racine obéit à l'équation : sommer la relation donne une récurrence sur les $s_k$.}
\begin{questions}
  \item $s_1=2$ et $s_2=-2$ (exercice 5).
  \item En sommant $x_i^{3}=2x_i^{2}-3x_i+4$ sur les trois racines :
    $s_3=2s_2-3s_1+12=-4-6+12=\res{2}$.
  \item En multipliant l'équation par $x_i^{k}$ puis en sommant :
    $s_{k+3}=2s_{k+2}-3s_{k+1}+4s_k$. Ainsi
    $s_4=2s_3-3s_2+4s_1=4+6+8=\res{18}$.
\end{questions}
\end{corrige}
\end{blocLivre}

\begin{blocLivre}{cor-10-21}
\begin{corrige}{21}{Racines rationnelles}
\demarche{Chasser les dénominateurs, puis faire travailler Gauss.}
\begin{questions}
  \item De $a_n p^{n}+a_{n-1}p^{n-1}q+\dots+a_0q^{n}=0$ : $p$ divise
    $a_0q^{n}$ et $p\wedge q=1$, donc (Gauss itéré) $p\mid a_0$ ;
    symétriquement $q\mid a_n$.
  \item $(\sqrt2+\sqrt3)^{2}=5+2\sqrt6$, donc
    $\bigl((\sqrt2+\sqrt3)^{2}-5\bigr)^{2}=24$ : le réel $\sqrt2+\sqrt3$
    annule $\res{X^{4}-10X^{2}+1}$.
  \item Les racines rationnelles éventuelles vérifient $p\mid1$ et
    $q\mid1$ : seules candidates $\pm1$, qui ne sont pas racines
    ($1-10+1=-8$). Donc toutes les racines sont irrationnelles, et en
    particulier $\res{\sqrt2+\sqrt3\notin\Q}$ — là où le chapitre 5
    passait par $\sqrt6$, l'arithmétique des polynômes conclut d'un coup.
\end{questions}
\end{corrige}
\end{blocLivre}

\begin{blocLivre}{cor-10-22}
\begin{corrige}{22}{L'erreur d'interpolation}
\demarche{Une fonction auxiliaire à quatre zéros, et Rolle en cascade.}
\begin{questions}
  \item $g$ s'annule en $x_1$, $x_2$, $x_3$ (où $f=P$ et le produit est
    nul) et en $x$ (par choix de $K$) : au moins \res{\text{quatre zéros}}.
  \item Rolle entre zéros consécutifs : $g'$ a au moins trois zéros, puis
    $g''$ deux, puis $g'''$ un — soit $\xi$ avec $g'''(\xi)=0$. Or
    $P'''=0$ (degré $\leq2$) et la dérivée troisième du produit unitaire
    de degré $3$ vaut $3!=6$ : $g'''=f'''-6K$, d'où
    $K=\dfrac{f'''(\xi)}{6}$.
  \item En majorant $|f'''(\xi)|$ par $M_3$ :
    \[
      |f(x)-P(x)|\leq\frac{M_3}{6}\,|(x-x_1)(x-x_2)(x-x_3)| :
    \]
    l'erreur est contrôlée par la dérivée troisième et la position des
    nœuds — le point de départ de la théorie de l'approximation.
\end{questions}
\end{corrige}
\end{blocLivre}

\begin{blocLivre}{cor-10-23}
\begin{corrige}{23}{Une primitive par les éléments simples}
\demarche{Pôle réel simple, puis trinôme sans racine réelle : logarithmes et arctangente.}
\begin{questions}
  \item $X^{3}+1=(X+1)(X^{2}-X+1)$. Le coefficient du pôle $-1$ vaut
    $\dfrac{1}{3}$ (évaluation de $\frac{1}{X^{2}-X+1}$ en $-1$), et
    l'identification donne
    \[
      \res{\frac{1}{X^{3}+1}=\frac{1}{3}\cdot\frac{1}{X+1}
      +\frac{-X+2}{3\,(X^{2}-X+1)}}.
    \]
  \item On écrit $\dfrac{-X+2}{X^{2}-X+1}
    =-\dfrac12\cdot\dfrac{2X-1}{X^{2}-X+1}
    +\dfrac32\cdot\dfrac{1}{(X-\frac12)^{2}+\frac34}$, d'où la primitive
    \[
      \res{\frac13\ln(x+1)-\frac16\ln(x^{2}-x+1)
      +\frac{1}{\sqrt3}\arctan\frac{2x-1}{\sqrt3}}+C.
    \]
  \item La dérivation redonne
    $\dfrac{1}{3(x+1)}+\dfrac{-x+2}{3(x^{2}-x+1)}=\dfrac{1}{x^{3}+1}$ —
    contrôle conforme (et vérifié numériquement).
\end{questions}
\end{corrige}
\end{blocLivre}

\begin{blocLivre}{cor-10-24}
\begin{corrige}{24}{Polynômes positifs et sommes de deux carrés}
\demarche{Factoriser sur $\R$, écrire chaque brique comme somme de deux carrés, multiplier avec l'identité des modules.}
\begin{questions}
  \item Développement direct — c'est l'identité
    $|zw|^{2}=|z|^{2}|w|^{2}$ des chapitres 3 et 9, lue sur les
    polynômes.
  \item Au voisinage d'une racine réelle de multiplicité impaire, $P$
    changerait de signe : exclu par $P\geq0$. Et le coefficient dominant
    est positif (limite en $+\infty$).
  \item La factorisation réelle s'écrit donc
    $P=c\,\prod(X-a_i)^{2m_i}\prod(X^{2}+p_jX+q_j)$ avec $c>0$ et
    $\Delta_j<0$ ; chaque trinôme s'écrit
    $\Bigl(X+\frac{p_j}{2}\Bigr)^{2}+\Bigl(\sqrt{q_j-\frac{p_j^{2}}{4}}\Bigr)^{2}$,
    somme de deux carrés, et le préfacteur
    $c\prod(X-a_i)^{2m_i}=\bigl(\sqrt c\prod(X-a_i)^{m_i}\bigr)^{2}+0^{2}$
    aussi.
  \item L'identité de la question 1 montre qu'un produit de sommes de
    deux carrés est une somme de deux carrés : de proche en proche,
    $P=A^{2}+B^{2}$. Exemple :
    $X^{4}+3X^{2}+2=(X^{2}+1)(X^{2}+2)$, et l'identité avec
    $(A,B,C,D)=(X,1,X,\sqrt2)$ donne
    \[
      \res{X^{4}+3X^{2}+2=\bigl(X^{2}-\sqrt2\bigr)^{2}
      +\bigl(1+\sqrt2\bigr)^{2}X^{2}}.
    \]
\end{questions}
\end{corrige}
\end{blocLivre}

% !TEX root = ../../main.tex
%% Corrigé du défi du chapitre 10 — vérifié par verifications/defis_06_10.py.

\begin{blocLivre}{cor-10-25a}
\begin{corrige}{25}{Le théorème de Gauss-Lucas}
\demarche{La dérivée logarithmique transforme « $P'(z)=0$ » en une somme
nulle ; en conjuguant, cette somme devient une relation barycentrique à
poids positifs.}
\partieExo{A}{La dérivée logarithmique}
\begin{questions}
  \item $P=\lambda\prod_k(X-a_k)^{m_k}$, et la dérivée d'un produit donne
    $P'=\lambda\sum_k m_k(X-a_k)^{m_k-1}\prod_{j\neq k}(X-a_j)^{m_j}$. En
    divisant par $P$ : $\res{\dfrac{P'}{P}=\sum_{k=1}^{p}\dfrac{m_k}{X-a_k}}$.
  \item En $z$ (qui n'annule pas $P$) :
    $0=\dfrac{P'(z)}{P(z)}=\sum_k\dfrac{m_k}{z-a_k}$. En conjuguant, et
    puisque $\dfrac{1}{\overline{z-a_k}}=\dfrac{z-a_k}{|z-a_k|^{2}}$ :
    $\res{\sum_k m_k\,\dfrac{z-a_k}{|z-a_k|^{2}}=0}$.
  \item Posons $w_k=\dfrac{m_k}{|z-a_k|^{2}}>0$ : la relation s'écrit
    $\bigl(\sum_kw_k\bigr)z=\sum_kw_ka_k$, donc $z=\sum_k\theta_ka_k$ avec
    $\theta_k=\dfrac{w_k}{\sum_jw_j}\geqslant0$ et $\sum\theta_k=1$. Si $z$
    est aussi racine de $P$, c'est l'un des $a_k$ ($\theta_k=1$). Toute
    racine de $P'$ est donc dans l'\res{\text{enveloppe convexe}} des
    racines de $P$.
\end{questions}
\partieExo{B}{Exemples}
\begin{questions}[resume]
  \item $P'=nX^{n-1}$ : seule racine $\res{0}$ (multiplicité $n-1$),
    centre du polygone régulier formé par les racines $n$-ièmes de
    l'unité — leur isobarycentre.
  \item $P=X^{3}-2\mathrm{i}X^{2}-X+2\mathrm{i}$, donc
    $P'=3X^{2}-4\mathrm{i}X-1$, de discriminant $-16+12=-4=(2\mathrm{i})^{2}$ :
    racines $\dfrac{4\mathrm{i}\pm2\mathrm{i}}{6}$, soit $\res{\mathrm{i}}$
    et $\res{\tfrac{\mathrm{i}}{3}}$. Et
    \[
    \mathrm{i}=\tfrac14\cdot1+\tfrac14\cdot(-1)+\tfrac12\cdot2\mathrm{i},
    \qquad
    \tfrac{\mathrm{i}}{3}=\tfrac{5}{12}\cdot1+\tfrac{5}{12}\cdot(-1)
    +\tfrac16\cdot2\mathrm{i} :
    \]
    deux barycentres à poids positifs, conformément au théorème.
\end{questions}
\begin{center}
\begin{tikzpicture}[scale=1.15,line cap=round,line join=round]
  \draw[otMuted,line width=0.4pt,-{Stealth[length=1.5mm]}] (-1.6,0) -- (1.7,0);
  \draw[otMuted,line width=0.4pt,-{Stealth[length=1.5mm]}] (0,-0.3) -- (0,2.4);
  \fill[otNavy!10] (1,0) -- (0,2) -- (-1,0) -- cycle;
  \draw[otNavy,line width=0.8pt] (1,0) -- (0,2) -- (-1,0) -- cycle;
  \foreach \p/\n/\a in {(1,0)/1/north,(-1,0)/-1/north,(0,2)/2\mathrm{i}/west}{%
    \fill[otNavy] \p circle (2pt);
    \node[font=\scriptsize,text=otNavy,anchor=\a] at \p {$\n$};}
  \filldraw[fill=otGreen,draw=otGreenDark,line width=0.5pt] (0,1) circle (1.9pt);
  \filldraw[fill=otGreen,draw=otGreenDark,line width=0.5pt] (0,0.3333) circle (1.9pt);
  \node[font=\scriptsize,text=otGreenDark,anchor=west] at (0.08,1) {$\mathrm{i}$};
  \node[font=\scriptsize,text=otGreenDark,anchor=west] at (0.08,0.36) {$\mathrm{i}/3$};
\end{tikzpicture}\\[0.5mm]
\legende{Les racines $\mathrm{i}$ et $\mathrm{i}/3$ de $P'$ sont dans le
triangle des racines de $P$.}
\end{center}
\end{corrige}
\end{blocLivre}

\begin{blocLivre}{cor-10-25b}
\begin{corrigeSuite}
\begin{questions}[start=6]
  \item Entre deux racines consécutives $a_k<a_{k+1}$, Rolle fournit une
    racine de $P'$ dans $\intOO{a_k}{a_{k+1}}$ : $p-1$ racines. Chaque
    $a_k$ est par ailleurs racine de $P'$ de multiplicité $m_k-1$. Au
    total, $(p-1)+\sum_k(m_k-1)=n-1=\deg P'$ racines dans
    $\intFF{a_1}{a_p}$ : \res{P'\text{ est scindé sur }\R}, et ses racines
    sont bien dans l'enveloppe convexe (le segment $\intFF{a_1}{a_p}$).
\end{questions}
\partieExo{C}{Conséquences}
\begin{questions}[start=7]
  \item Si $P=\lambda\bigl(X^{n}-\sigma_1X^{n-1}+\cdots\bigr)$, où
    $\sigma_1$ est la somme des racines, alors
    $P'=\lambda\bigl(nX^{n-1}-(n-1)\sigma_1X^{n-2}+\cdots\bigr)$, dont la
    somme des racines vaut $\frac{(n-1)\sigma_1}{n}$ (Viète). Les moyennes
    valent toutes deux $\res{\dfrac{\sigma_1}{n}}$.
  \item Un disque est convexe : il contient tous les barycentres à poids
    positifs de ses points. Par Gauss-Lucas, les racines de $P'$ sont dans
    $D$ ; en réappliquant le théorème à $P'$, celles de $P''$ aussi, et
    ainsi de suite. Enfin
    $P^{(n-1)}=\lambda\bigl(n!\,X-(n-1)!\,\sigma_1\bigr)$ a pour unique
    racine $\res{\dfrac{\sigma_1}{n}}$, l'isobarycentre des racines de $P$ :
    les dérivées successives « contractent » les racines vers leur centre
    de gravité.
\end{questions}
\end{corrigeSuite}
\end{blocLivre}

